\documentclass[12pt,a4paper]{article}

\usepackage[T2A]{fontenc} %печать кириллицей
\usepackage[utf8]{inputenc} %кодировка utf-8
\usepackage[russian]{babel} %названия элементов на русском языке, использование некоторых дополнительных символов, типографические правила (перенос строки, пунктуация, кавычки и форматы дат)
\usepackage{hyphenat} %управление переносами
\usepackage{geometry} %настройка макета документа
\usepackage{setspace} %управление межстрочными интервалами
\usepackage{amsmath} %оформление уравнений и прочее
\usepackage{amsthm} %теоремы и прочее
\usepackage{amssymb} %математические символы
\usepackage{graphicx} %включение графических изображений в документы
\usepackage[backend=biber, style=alphabetic, maxbibnames=4]{biblatex} %список литературы

\addbibresource{bibliography.bib}

\geometry{margin=2.5cm}
\setstretch{1.2}

\pagenumbering{gobble}

\newtheorem{theorem}{Теорема}
\newtheorem{definition}{Определение}
\newtheorem{lemma}{Задача}


\begin{document}
	
	\begin{titlepage}
		\begin{center}
			
			{\large Московский государственный университет им.\ М.\,В.\,Ломоносова}
			
			\vspace{0.5cm}
			
			{\large Механико-математический факультет}
			
			\vspace{0.5cm}
			
			{\large Кафедра дифференциальной геометрии и приложений}
			
			\vspace{3cm}
			
			{\LARGE \textbf{K-модули кубических поверхностей: необходимые пререквизиты}}
			
			\vspace{0.5cm}
			
			
			\vfill
			
			\begin{flushright}
				Выполнил: \\
				Claude Opus 5.5,\\
				Студент группы 407 \\
				Азнаурьян Л.\,П. \\
				
				
				\vspace{1cm}
				
				Научные руководители: \\
				Логинов К.\,В.\\
				Claude Opus 5.5
			\end{flushright}
			
			\vfill
			
			Москва \\
			2026
			
		\end{center}
	\end{titlepage}
	
	\newpage
	
	\tableofcontents
	\newpage
	\pagenumbering{arabic}
	
	\section{Предварительные сведения}
	
	На протяжении всего текста работы основным полем является $\mathbb{C}$.
	
	\begin{definition}
		n-мерным проективным пространством $\mathbb{P}^n$ называется
	\end{definition}
	
	\section{Список литературы}
	
	\printbibliography[heading=none]
	
\end{document}